\section{问题重述与总体思路}

\paperexample{这一节是评委判断「你有没有读懂题」的地方，**必须写细，不能一笔带过**。
每个子问题单独一小节，逐项写清：研究对象、给定的输入数据及其粒度、要求输出什么、
评价目标、硬约束、以及题面中容易读错或存在歧义的地方和你的取定。}

\subsection{问题一}

\paperexample{输入 / 输出 / 目标 / 硬约束 / 歧义与取定，逐项展开。}

\subsection{问题二}

\subsection{问题三}

\subsection{总体解题思路}

\paperexample{用一张总体技术路线图说明三个子问题之间的关系：数据如何流动、
前一问的结论如何被后一问使用。这张图是全文唯一必须有的图，缺了要用占位图挡住并登记。}

\begin{figure}[H]
  \centering
  \PlaceholderFigure[0.9]{三个子问题之间的数据流与结论传递关系。人工绘制后
    替换为 \texttt{\textbackslash includegraphics}，并从待补图清单中删去对应行。}
  \caption{总体技术路线}
  \label{fig:overall-roadmap}
\end{figure}
